\section{Mål}
\label{sec:maal}

\subsection{Effektmål}

VideoAPI skal bidra til økt fleksibilitet ved integrasjon av nye videomotorer med minimale endringer i eksisterende løsninger. Løsningen skal redusere teknologisk innlåsing ved å gjøre det enklere å bytte leverandør basert på pris, funksjonalitet eller driftsbehov.

\subsection{Resultatmål}    

Prosjektet resulterte i følgende leveranser:

\begin{enumerate}
    \item Et REST-API med plugin-arkitektur, OAuth2/JWT-autentisering,
          health-endepunkter, strukturert logging og Prometheus-monitorering.
    \item Provider-adaptere for AntMedia Server og NHN Video.
    \item En enkel demo-frontend for å demonstrere API-et i praksis.
    \item OpenAPI/Swagger-dokumentasjon beregnet på eksterne leverandører.
    \item README, deployment-guide og utviklerdokumentasjon for videre arbeid.
    \item En containerisert løsning med Docker og CI/CD via GitHub Actions.
    \item Denne bachelorrapporten.
\end{enumerate}

\subsection{Læringsmål}

Teamet hadde ingen tidligere erfaring med C\# eller .NET. De første sprintene ble derfor benyttet til å tilegne seg språket og økosystemet parallelt med utvikling. Gjennom prosjektet opparbeidet teamet praktisk erfaring med .NET 10, ASP.NET Core, REST-API-design, plugin-arkitektur, Docker, GitHub Actions og Prometheus/Grafana. Samarbeidet med HDO som ekstern oppdragsgiver ga også erfaring med kravhåndtering, tekniske avklaringer og arbeid under leveransekrav i et reelt prosjekt.